\chapter[Sermon 7. Of the Epistles Revealed Out of the Throne of God from Christ by an Angel, and Received and Sent of John. Where Also a Part of the Epistle to the Ephesians Is Expounded.]{Of the Epistles Revealed Out of the Throne of God from Christ by an Angel, and Received and Sent of John. Where Also a Part of the Epistle to the Ephesians Is Expounded.}
\label{sermon:007}
\markright{Sermon 7. Of the Epistles Revealed Out of the Throne of God from ...}

To the messenger of the congregation of Ephesus write: These things says he who holds the seven stars in his right hand and walks among the seven golden lampstands. I know your works, your labor, and your patience, and how you cannot bear those who are evil. You examine those who claim to be apostles but are not, and you have found them liars. You have endured, showing patience. For my name’s sake you have labored and have not fainted. Nevertheless, I have something against you, because you have left your first love.

Your kindness has observed a certain image of the Lord Christ, sitting on the right hand of the Father in glory; yet so, that in no wise does he forsake or neglect his Church. Now follows a more full and plain explanation of how our Savior Christ in Heaven executes the office of the high Bishop, and teaches the whole church by his ministers, rebukes, comforts, and retains it in its duty; finally, turns away all things hurtful, and advances it to greater things. For here follow seven Epistles, to the seven Congregations: that is to say, unto all the churches in the [illegible] whole world. For this most ample and wholesome doctrine [illegible] must not be restrained to a few, since Christ is Bishop universal.

But great is the authority of these epistles. For they are revealed from the throne of God, by the Son of God speaking by an angel, which prescribes what is to be written in those epistles. John receives and writes the same, through Christ’s commandment, and sends them to the seven congregations. And verily they pertain no less unto us than if now the bearer entering into the church should deliver these letters unto us.

Moreover, in these seven Churches is figured unto us the nature, manners, vices, medicines, rebukes, praises of all Churches in all times, and whatsoever is wont to happen about them. Then by examples of most excellent means mixed, of hypocritical also and wicked. And these our Lord doth evidently instruct, reprove, rebuke, and blame, praise, correct, move, exhort, comfort, the same he threatens, and promises them also joyful things, \&c. This is no light or common example, but of the Son of God, the high and most blessed Bishop: teaching us how we should deal with congregations after the capacity and disposition of every one.

And not without cause he chooses seven of the most noble cities of Asia. It is certain that Asia was among the first regions inhabited, and from there people were dispersed into other parts of the world. It is also certain that the Devil set up his throne in Asia, and there reigned in men through idolatry, murder, ambition, avarice, uncleanness, and base pleasures. For the proverb says, “The laughter of Jove.” It is known what the apostle wrote to the Ephesians in the fourth and fifth chapters. Therefore, our Savior Christ would overthrow that throne of the Devil and set up his throne of righteousness and holiness. Therefore, he goes first and chiefly to those of Asia, that by their example the whole world might be corrected and amended.

\index{Rome}And amongst other cities of Asia and Ionia, Ephesus was most famous, called in the old time the light of Asia. And amongst the twelve cities of Ionia, it was accounted the principal port. Neither was there any other Church richer or more beautiful anywhere in Asia than the temple of Diana at Ephesus. It stood in the midst of the city, a great wonder of Greek magnificence, as writes Pliny. This temple is said to have been built in two hundred and twenty years by a king of all Asia, and set in a marshy ground so that it would not feel earthquakes or the opening of the earth. The length of the temple was 425 feet, the breadth 220 feet. It had 107 pillars, dedicated by so many kings. Seek the rest out of the summary by the famous Dr. Joachim Vadian. The apostle Paul first illuminated this city with the name of the Gospel. His epistle to the same city remains, and a plentiful story in the Acts of the Apostles. After Paul was executed, John went to Ephesus, and from thence preached to all Asia; from thence also he was brought to Rome to the emperor Domitian. To Ephesus he returned after his exile, and there at length, as ecclesiastical histories testify, he slept in the Lord.

And before every epistle, especially that of the Ephesians, is placed a command, write. This command gives authority to the writing, so that we may not ask, whether this writing ought to be credited? and why it should be believed? For here is the express commandment of God, and the divine authority, concerning which to inquire curiously is thought not without cause unlawful. Moses wrote by the commandment of God. And by the same commandment of God wrote also the prophets and Apostles. Which writings are not believed to be authentic? Certainly John said truly and wittily: he who knows God hears us; he who is not of God hears us not. 1 John 5. Curious questions cease, where the mind of the godly, or of any poor sheep, knows the voice of his Lord and Shepherd.

And let no man think that this letter, being written to one angel—that is, a bishop or pastor—appertains nothing to the Church. For to the end of the letter is added an acclamation: “He that hath an ear let him hear what the Spirit says to the congregations.” Therefore, the pastor is named, but the sheep are not excluded. All ranks and states in the church know what is said to them. Ignorance says: “That which is written to the Romans concerns me nothing.” Yet nevertheless it is addressed to the angel, to the intent that pastors may be admonished concerning the state of the Church.

\index{Repentance}The argument of the first epistle is this. Christ declares that he rules over his church, that he takes charge and government of the same. Some things he praises, and some he blames. In the meantime, he exhorts to repentance, threatening grievous things and promising most joyful ones. And also, he applies this epistle to all churches and communicates it to all congregations in the whole world. But the epistle is exhortative, for it instructs the churches, exhorts, and directs.

And first, he shows who he is, from whom the epistle proceeds, that he may give authority thereunto; and may declare also that he is the head of his church, the Bishop, Duke, and governor. That part is taken from the image set forth in the first chapter. And follows the prophetic manner of speaking: This says he who holds the seven stars in his right hand. For the Prophets say likewise: “Thus says the Lord God of Israel,” “Thus says the Lord of hosts.” “Thus says the Lord,” who brought you out of Egypt, \&c. And two special things he repeats from the former description, whereby he will be known and do us understand how he being Lord and Bishop rules and maintains in his church. First, he affirms that he holds in his hand the seven stars. The hand is a token of working, protection, or deliverance. The stars we have heard are the ministers, and the ministry of the word, or the church. Therefore Christ holds the ministry in the church, and the ministers work the salvation of the faithful. After, he affirms that he walks, not sleeping or doing nothing, in the midst of seven golden candlesticks. In the midst said, to the end we should understand, that he gives himself indifferently to all men, and rules over all with like power and government. Fulwel wrote hereof. What, says he, is to walk or to be in the midst of congregations, but to assist them, keep, instruct, help them, and by all means to watch over them? For which cause he says also in the last of Saint Matthew: “Behold, I am with you always, even unto the end of the world.” Hereof you have a most excellent figure in the law: wherein amongst other things which pertained to the ministry of the high priest, he had charge of oil and of seven candles; for those must he snuff and pour in oil, when it wanted. So Christ, the high and true Bishop, has the charge of the seven candles, that is to say, of all congregations; and is careful that they want not that oil, which is mentioned in the 44th Psalm. He watches that they want not the fire and light of the truth. Finally, he snuffs and purges by faith, whatever thing so ever has need to be purged in them. Thus far he. Which things when they hear, which make the Bishop of Rome head of the church, it is marvelous if by and by they understand not their folly and madness. Here the Lord adds also, that he knows the works, to wit, all both good and evil, as well of the Bishop as of his church. For the Lord knows all things, and is head Bishop of the Catholic or universal Church, which also remembers the thoughts of all men in the world at one instant; who sees what is done, and what is not done, and what things are needful, nothing escapes him. And such in deed ought he to be, that is head universal of his church. And this sentence is repeated, “I know thy works,” in the beginning of every epistle. And verily it is full of comfort, when we hear that Christ knows all our doings. For we believe also that he has a faithful care of all our matters.

Now this great Bishop commends certain things within the congregation of Ephesus. For good works in deed are approved by Christ, and he praises them, so that he might encourage those who run in his way. First, he acknowledges the labor and patience of both the Bishop and the Church. Labor encompasses thought and care in the way of God, mortifying the flesh, the study of good works, but chiefly the cross and persecution, which history testifies to have been extreme and cruel in the time of Domitian. And unless the persecuted have patience, they cannot endure the labor. Holy patience keeps us in work and holy labor.

But lest patience should be stretched to those things wherein impatience is considered praiseworthy, he adds the second point that he praises in them, that you cannot bear evil men. And by these evil men he means not weaklings, or those who err without malice: but the prophet David also says, Psalm 119, “I have hated the wicked; your law I have loved.” What we should do with the weak in the faith, or with those who err through ignorance rather than obstinate stubbornness, the Apostle has taught us in Romans 14. The example of our Savior has also taught, bringing again that strayed sheep upon his shoulder into the sheepfold. Therefore, the Lord speaks here of the obstinate, of the deceivers who delight to err themselves and to draw others with them into errors; no Christian patience bids us to bear with such men.

\index{Jerome, Saint}\index{Irenaeus}\index{Eusebius of Caesarea}\index{Justification}And in the words following he declares of what sort those evil men were. And you have examined those who say they are Apostles, but are not, and have found them liars. Observe, he speaks of the false apostles, of whom in John’s time there was exceedingly great abundance. For they were Zarbians, mixing law with grace, and attributing justification to the law and to our own righteousness. Which the holy and great council at Jerusalem condemned, as it appears in the fifteenth chapter of the Acts of the Apostles. Such a false apostle was Hebion. Eusebius mentions in the book of Ecclesiastical History the twenty-seventh chapter. Here was added Cerinthus, that heretic, not apostle. There were more also, of whom some denied the humanity of Christ, some his deity. Against whom John wrote in his Gospel and in his Epistle, and Irenaeus in the first book against heretics. These the Lord denies to be apostles, or apostolic men, which the apostles have also denied, Acts 15. And therefore the apostle Jerome in his canonical epistle says, “Who is a liar, but he who denies Jesus to be Christ?” But if such trouble were in the churches while the apostles were yet living, and there were then so many deceivers, what marvel is it, though in the dregs of the world, that is, in this our last time, there are not a few sounding where they are now that twisted dissensions and troubles in the defense of their error? The Gospelers themselves say they are at variance. God is God of concord; how then should I believe that God is among those who cause strife? So might the Sophists also have reasoned in the apostles’ time.

And here we have a fitting way, in what manner the churches should proceed, while troublesome persons, like false Apostles, attempt to divide the Church further. For such ringleaders must be tried and examined: and they must be tried according to the Christian belief and doctrine of the Apostles, and inquiry must be made whether they are Apostles and true men, or false Apostles and liars. When we have found them to be false Apostles and liars, and they persist stubbornly in their wickedness, they are not to be tolerated; as indeed the Ephesians did not deem safe to bear with such deceivers. And we must know that pastors ought to act one way, the Christian magistrate another way, and the people a third way, not to tolerate open heretics. For the pastor not only does not tolerate them, by being watchful and wary of those wolves, but assails them with wholesome doctrine, and rejects them from the flocks of Christ. But the magistrate, because he is a Christian magistrate, and by his duty also, not only as a private person but also as a magistrate, ought to serve Christ; he ought also with the sword of justice to drive away poison from the church, and to punish manifest blasphemies. And the people are commanded neither to hear them, nor receive them, nor to have anything to do with heretics, and so not to tolerate them. They may therefore be ashamed of their faultiness and the persistence of their perverse patience, which think it no shame to maintain heretics, and to flatter the manifest enemies of Christ and the Church. Psalm 15. He is praised who does not value the wicked; that is to say, in whose sight the wicked man is vile. Therefore, he is rightly blamed, whoever flatters the ungodly. And the hatred in deed is rather against wickedness than against the person of the wicked, which itself is commanded to be loved. The Devil at this day raises up the old heresies of Hiero[?]m, Ireney[?], and others in Ter[illegible]paniarde, and in the Anabaptists, Libertines, and other monsters, so that the thing itself and the danger there[illegible] commands us to watch and to drive away the most ravenous wolves from the holy Church of Christ, which nevertheless set forth nothing more than patience and charity, for this intent verily that they might be spared, and might unpunished teach what they list against Christ, and work against his church, yea tear it in pieces with their wicked teeth.

But when these evil men are not tolerated, but oppose those who deceive and are deceived, a great conflict arises, involving labor, thought, anxiety, watchfulness, and injuries to be suffered for the name of Christ and the defense of the truth. For unless we are diligent and patient here, the deceivers will prevail. In this, the church of the Ephesians behaved admirably, so that the Lord now commends exceedingly the courage, patience, and constancy of the pastor and his church. For these things should not be understood to refer to that patience whereby evil men are tolerated and permitted to proceed in their malice and deceitfulness. Otherwise, this passage would contradict what preceded it. This is what the learned interpreter seems to have intended to avoid. For he reads, “You have patience, and have suffered,” where the Greek says, “You have suffered and have patience.” He altered the order, and would not place “you have suffered” before “you have patience,” lest anyone misunderstand that they had tolerated the false apostles. Instead, he placed “patience” first and “suffered” afterward, so that we might understand that they did not tolerate evil men, but the evil wrought by evil men. Thus, they labored with patience for Christ’s name, namely, to maintain it against harmful heresies. And he adds, “You have not fainted, though weary and broken with labors.” For we are taught to overcome through patient constancy, which is rightly called the fulfillment of every good work.

All and every of these things we must apply to ourselves and understand how we may also, this day, please Christ our Redeemer, King and Bishop. Truly, we walk in the same steps wherein we see the congregation of the Ephesians to have walked.

What thing did he condemn in that church? That they had abandoned their first love. When they first received the Gospel by Paul, and afterward by John and other godly men, there was seen a great fervor in the words and deeds of the faithful; which thing may be gathered both from the acts of the Apostles and also from Paul’s epistle to the Ephesians. They loved God and their neighbors with a most fervent zeal. They burned in reforming of manners. But in process of time, this heat was cooled, and they grew colder in true godliness. This great mischief he rebukes in them, and then desires to have it redressed. And here let us note that not only revolting from idolatry and other great crimes are imputed to the church, but also if we slack in any thing in holy zeal. So that hereof we may learn how holy and blameless we ought to be before God. Doubtless we cannot here excuse ourselves before the divine majesty, who has been for thirty years past more fervent in this congregation than we are today. \&c. Our Lord God enlighten our minds, that we may please him. To whom be glory.
